% Source: https://artofproblemsolving.com/community/c3h1928316_the_boxes
% Stable identity: aops-topic-1928316-q01
Every label is wrong, so the contents are a rearrangement of the labels with no box in its own place. There are only two such arrangements:
\begin{itemize}
\item Arrangement 1: ``APPLES'' holds oranges, ``ORANGES'' holds the mixture, ``APPLES AND ORANGES'' holds apples.
\item Arrangement 2: ``APPLES'' holds the mixture, ``ORANGES'' holds apples, ``APPLES AND ORANGES'' holds oranges.
\end{itemize}
(The box labelled ``APPLES AND ORANGES'' holds apples or oranges, and once its contents are known, the other two boxes are forced.)

\textbf{Choose the box labelled ``APPLES AND ORANGES''.} It contains only one kind of fruit, so an apple means Arrangement 1 and an orange means Arrangement 2. In either case all three labels are now known: for example, after an apple, the box labelled ``ORANGES'' cannot hold oranges (wrong label) or apples (already placed), so it holds the mixture, and the box labelled ``APPLES'' holds oranges.

\textbf{The other choices can fail.} From the box labelled ``APPLES'' you may draw an orange in both arrangements (it holds oranges or the mixture). From the box labelled ``ORANGES'' you may draw an apple in both arrangements (it holds the mixture or apples). In those cases the fruit does not tell the two arrangements apart.
